% ============================================================
% Secao 10: Continual Learning
% ============================================================

\section{Continual Learning}
\label{sec:continual}

O \aletheion{} integra dois mecanismos complementares de
\textit{Continual Learning} (CL) para preservar conhecimento ao longo
de m\'ultiplas fases de treinamento: \textit{Elastic Weight Consolidation}
(EWC) e \textit{Experience Replay}.

\subsection{Motiva\c{c}\~ao}

O \textit{catastrophic forgetting} \cite{mccloskey1989catastrophic} \'e
o fen\^omeno em que redes neurais perdem conhecimento de tarefas
anteriores ao treinar em novas tarefas. Em um LLM com sistema
epist\^emico, isso \'e particularmente problem\'atico:

\begin{itemize}
    \item O manifold 5D pode ``esquecer'' sua geometria calibrada
    \item O $\tau^2$ pode descalibrar-se
    \item Os sub-heads epist\^emicos podem perder sinais aprendidos
\end{itemize}

\subsection{EWC: Elastic Weight Consolidation}
\label{sec:ewc}

EWC \cite{kirkpatrick2017ewc} protege par\^ametros importantes adicionando
uma penalidade quadr\'atica ponderada pela \textit{Fisher Information Matrix}
(FIM) diagonal.

\subsubsection{Fisher Information Matrix}

Ap\'os cada fase de treinamento, estimamos a FIM diagonal:

\begin{equation}
    F_i = \E\left[\left(\frac{\partial \calL}{\partial \theta_i}\right)^2\right]
    \approx \frac{1}{N} \sum_{n=1}^{N} \left(\frac{\partial \calL(\bx^{(n)})}{\partial \theta_i}\right)^2
    \label{eq:fisher}
\end{equation}

onde $N = 256$ amostras do dataset (configur\'avel via
\texttt{ewc\_fisher\_samples}).

$F_i$ mede a \textbf{import\^ancia} do par\^ametro $\theta_i$: valores altos
indicam que o par\^ametro \'e cr\'itico para a tarefa atual.

\subsubsection{Penalidade EWC}

Na fase seguinte, a loss total inclui:

\begin{equation}
    \calL_{\text{EWC}} = \frac{\lambda_{\text{EWC}}}{2} \sum_{i} F_i \cdot (\theta_i - \theta_i^*)^2
    \label{eq:ewc_loss}
\end{equation}

onde $\theta_i^*$ s\~ao os par\^ametros ap\'os a fase anterior e
$\lambda_{\text{EWC}} = 100.0$ controla a rigidez.

\subsubsection{EWC Online}

Para m\'ultiplas fases, utilizamos EWC online
\cite{schwarz2018progress} com acumula\c{c}\~ao exponencial:

\begin{align}
    \bar{F}_i^{(k)} &= \gamma \cdot \bar{F}_i^{(k-1)} + F_i^{(k)} \label{eq:ewc_online} \\
    \theta_i^{*(k)} &= \theta_i^{(k)}
\end{align}

com $\gamma = 0.9$ (decay). Isso evita o crescimento linear de termos
que ocorre no EWC multi-tarefa ingenuamente.

\subsubsection{Integra\c{c}\~ao no Trainer}

A \Cref{eq:ewc_loss} \'e adicionada {\`a} loss total ap\'os o c\'alculo da
loss composta:

\begin{equation}
    \calL_{\text{final}} = \calL_{\text{total}} + \calL_{\text{EWC}}
\end{equation}

O m\'etodo \texttt{consolidate\_phase()} deve ser chamado ap\'os cada
fase:

\begin{lstlisting}[caption={Uso de EWC no treinamento.}]
# Fase 1: Treinamento base
trainer.train()
trainer.consolidate_phase()  # Computa Fisher

# Fase 2: Novo dataset (EWC protege fase 1)
trainer.train_loader = new_loader
trainer.train()
trainer.consolidate_phase()  # Acumula Fisher (online)

# Fase 3: Mais dados (EWC protege fases 1+2)
trainer.train_loader = another_loader
trainer.train()
\end{lstlisting}

\subsection{Experience Replay}
\label{sec:replay}

\textit{Experience Replay} complementa o EWC mantendo um buffer de
amostras anteriores que s\~ao misturadas ao batch atual.

\subsubsection{Reservoir Sampling}

O buffer utiliza reservoir sampling de Vitter \cite{vitter1985reservoir}
para manter uma amostra uniforme de tamanho fixo:

\begin{algorithm}[H]
\caption{Reservoir Sampling}
\begin{algorithmic}[1]
\REQUIRE Buffer $\calB$ de capacidade $M$, amostra $\bx_n$ (n-\'esima)
\IF{$|\calB| < M$}
    \STATE $\calB \leftarrow \calB \cup \{\bx_n\}$
\ELSE
    \STATE $j \sim \text{Uniform}(0, n)$
    \IF{$j < M$}
        \STATE $\calB[j] \leftarrow \bx_n$
    \ENDIF
\ENDIF
\end{algorithmic}
\end{algorithm}

Isso garante que cada amostra vista at\'e o momento $n$ tem probabilidade
$M/n$ de estar no buffer, sem necesitar armazenar todo o hist\'orico.

\subsubsection{Mix Batch}

Em cada step de treinamento, uma fra\c{c}\~ao $\rho = 0.1$ do batch
\'e substitu\'ida por amostras do buffer:

\begin{equation}
    \tilde{\calB}_t = (1 - \rho) \cdot \calB_t^{\text{novo}} \cup \rho \cdot \text{sample}(\calB_{\text{replay}})
\end{equation}

\subsubsection{Armazenamento CPU}

O buffer \'e mantido em mem\'oria CPU para n\~ao consumir VRAM:

\begin{equation}
    \text{mem}_{\text{buffer}} = M \cdot T \cdot 2 \text{ bytes} \approx 10000 \cdot 1024 \cdot 2 \approx 20 \text{ MB}
\end{equation}

Amostras s\~ao transferidas para GPU apenas quando selecionadas para o
\textit{mix batch}.

\subsection{Sinergia EWC + Replay}

Os dois mecanismos s\~ao complementares:

\begin{itemize}
    \item \textbf{EWC} protege no \textit{espa\c{c}o de par\^ametros}:
    impede que pesos importantes mudem
    \item \textbf{Replay} protege no \textit{espa\c{c}o de dados}:
    mant\'em exemplos do passado no treinamento
\end{itemize}

Juntos, reduzem o \textit{forgetting} mais do que qualquer um isoladamente
\cite{chaudhry2019continual}.

\subsection{Checkpoint com Estado CL}

O estado de CL \'e salvo junto ao checkpoint do modelo:

\begin{lstlisting}[caption={Conte\'udo do checkpoint.}]
checkpoint = {
    'model': model.state_dict(),
    'optimizer': optimizer.state_dict(),
    'scheduler': scheduler.state_dict(),
    'global_step': step,
    'ewc_state': ewc.state_dict_ewc(),   # Fisher + params ref
    'replay_state': replay.state_dict(),  # Buffer + contadores
}
\end{lstlisting}

Isso permite retomar treinamento com CL intacto.

\subsection{Configura\c{c}\~ao}

\begin{table}[H]
\centering
\caption{Par\^ametros de Continual Learning.}
\label{tab:cl_config}
\begin{tabular}{lrl}
\toprule
\textbf{Par\^ametro} & \textbf{Default} & \textbf{Descri\c{c}\~ao} \\
\midrule
\texttt{enable\_ewc} & true & Ativa EWC \\
\texttt{ewc\_lambda} & 100.0 & Peso da penalidade EWC \\
\texttt{ewc\_online} & true & Acumula\c{c}\~ao online de Fisher \\
\texttt{ewc\_gamma} & 0.9 & Decay da Fisher acumulada \\
\texttt{ewc\_fisher\_samples} & 256 & Amostras para estimar Fisher \\
\texttt{enable\_replay} & true & Ativa Experience Replay \\
\texttt{replay\_buffer\_size} & 10000 & Capacidade do buffer \\
\texttt{replay\_mix\_ratio} & 0.1 & Fra\c{c}\~ao de replay no batch \\
\bottomrule
\end{tabular}
\end{table}
